\documentclass[12pt]{article}
\usepackage[a4paper,margin=1in]{geometry}
\usepackage[utf8]{inputenc}
\usepackage[T1]{fontenc}
\usepackage{amsmath,amssymb}
\usepackage{textcomp}
\usepackage{parskip}
\begin{document}
\begin{center}
{\Large \textbf{2010 STEP I}}
\end{center}

\subsection*{Page 1}
\begin{flushleft}
\bigskip\textbf{Section A: Pure Mathematics}\par
\bigskip\textbf{Question 1.} Given that\par
2 2 2 2\par
5x + 2y 6xy + 4x 4y a (x y + 2) + b (cx + y) + d ,\par
- - \equiv -\par
find the values of the constants a, b, c and d.\par
Solve the simultaneous equations\par
2 2\par
\[5x + 2y 6xy + 4x 4y = 9 ,\]
- -\par
2 2\par
\[6x + 3y 8xy + 8x 8y = 14 .\]
- -\par
x a\par
x\par
\bigskip\textbf{Question 2.} The curve y = - e , where a and b are constants, has two stationary points. Show that\par
x b\par
³ ´\par
-\par
\[a b < 0 or a b > 4 .\]
- -\par
\textbf{(i)} Show that, in the case a = 0 and b = 1 , there is one stationary point on either side of\par
the curve's vertical asymptote, and sketch the curve.\par
\textbf{(ii)} Sketch the curve in the case a = 9 and b = 0 .\par
\bigskip\textbf{Question 3.} Show that\par
\[sin(x + y) sin(x y) = 2 cos x sin y\]
- -\par
and deduce that\par
\[sin A sin B = 2 cos 1 (A + B) sin 1 (A B) .\]
- 2 2 -\par
Show also that\par
\[cos A cos B = 2 sin 1 (A + B) sin 1 (A B) .\]
- - 2 2 -\par
The points P , Q, R and S have coordinates (a cos p, b sin p), (a cos q, b sin q), (a cos r, b sin r)\par
and (a cos s, b sin s) respectively, where 0 6 p < q < r < s < 2\pi, and a and b are positive.\par
Given that neither of the lines P Q and SR is vertical, show that these lines are parallel if and\par
only if\par
\[r + s p q = 2\pi .\]
- -\par
\[9465 Jun10\]
\end{flushleft}
\subsection*{Page 2}
\begin{flushleft}
\bigskip\textbf{Question 4.} Use the substitution x = , where t > 1, to show that, for x > 0,\par
\[t^2 1\]
-\par
\[dx = 2 ln \sqrt\{x\} + \sqrt\{x\} + 1 + c .\]
Z x (x + 1)\par
¡ ¢\par
p\par
1 1 t a\par
\[[Note: You may use without proof the result dt = ln - + constant. ]\]
\[Z t^2 a^2 2a ¯ t + a ¯\]
- ¯ ¯\par
¯ ¯\par
The section of the curve ¯ ¯\par
1 1\par
\[y =\]
\[\sqrt\{x\} - \sqrt\{x\} + 1\]
\[between x = 1 and x = 9 is rotated through 360 o about the x-axis. Show that the volume\]
8 16\par
\[enclosed is 2\pi ln 5 .\]
n\par
\bigskip\textbf{Question 5.} By considering the expansion of (1 + x) where n is a positive integer, or otherwise, show\par
that:\par
n n n n n\par
\textbf{(i)} + + + + = 2 ;\par
µ0¶ µ1¶ µ2¶ \cdots µn¶\par
n n n n n-1\par
\textbf{(ii)} + 2 + 3 + + n = n2 ;\par
µ1¶ µ2¶ µ3¶ \cdots µn¶\par
n 1 n 1 n 1 n 1 n+1\par
\textbf{(iii)} + + + + = 2 1 ;\par
µ0¶ 2 µ1¶ 3 µ2¶ \cdots n + 1 µn¶ n + 1 -\par
¡ ¢\par
n 2 n 2 n 2 n n-2\par
\textbf{(iv)} + 2 + 3 + + n = n (n + 1) 2 .\par
µ1¶ µ2¶ µ3¶ \cdots µn¶\par
x\par
\bigskip\textbf{Question 6.} Show that, if y = e , then\par
\[d2y dy\]
\[(x 1) x + y = 0 . ( )\]
\[- dx2 - dx ∗\]
x\par
In order to find other solutions of this differential equation, now let y = ue , where u is a\par
function of x. By substituting this into ( ), show that\par
∗\par
\[d2u du\]
\[(x 1) + (x 2) = 0 . ( )\]
\[- dx2 - dx ∗∗\]
du\par
\[By setting = v in ( ) and solving the resulting first order differential equation for v,\]
dx ∗∗\par
x\par
\[find u in terms of x. Hence show that y = Ax + Be satisfies ( ), where A and B are any\]
∗\par
constants.\par
\[9465 Jun10 [Turn over\]
\end{flushleft}
\subsection*{Page 3}
\begin{flushleft}
7 Relative to a fixed origin O, the points A and B have position vectors a and b, respectively.\par
(The points O, A and B are not collinear.) The point C has position vector c given by\par
\[c = αa + βb ,\]
\[where α and β are positive constants with α + β < 1 . The lines OA and BC meet at the\]
point P with position vector p and the lines OB and AC meet at the point Q with position\par
vector q. Show that\par
αa\par
\[p = ,\]
1 β\par
-\par
and write down q in terms of α, β and b.\par
Show further that the point R with position vector r given by\par
αa + βb\par
\[r = ,\]
α + β\par
lies on the lines OC and AB.\par
OQ OS\par
\[The lines OB and P R intersect at the point S. Prove that = .\]
BQ BS\par
\bigskip\textbf{Question 8.} (i) Suppose that a, b and c are integers that satisfy the equation\par
3 3 3\par
\[a + 3b = 9c .\]
Explain why a must be divisible by 3, and show further that both b and c must also be\par
\[divisible by 3. Hence show that the only integer solution is a = b = c = 0 .\]
\textbf{(ii)} Suppose that p, q and r are integers that satisfy the equation\par
4 4 4\par
\[p + 2q = 5r .\]
By considering the possible final digit of each term, or otherwise, show that p and q are\par
\[divisible by 5. Hence show that the only integer solution is p = q = r = 0 .\]
\[9465 Jun10\]
\end{flushleft}
\subsection*{Page 4}
\begin{flushleft}
\bigskip\textbf{Section B: Mechanics}\par
2a\par
2b\par
α\par
The diagram shows a uniform rectangular lamina with sides of lengths 2a and 2b leaning\par
against a rough vertical wall, with one corner resting on a rough horizontal plane. The plane\par
of the lamina is vertical and perpendicular to the wall, and one edge makes an angle of α with\par
the horizontal plane. Show that the centre of mass of the lamina is a distance a cos α + b sin α\par
from the wall.\par
The coefficients of friction at the two points of contact are each µ and the friction is limiting\par
at both contacts. Show that\par
\[a cos(2λ + α) = b sin α ,\]
\[where tan λ = µ.\]
\[\pi\]
\[Show also that if the lamina is square, then λ = α.\]
4 -\par
\bigskip\textbf{Question 10.} A particle P moves so that, at time t, its displacement r from a fixed origin is given by\par
t t\par
\[r = e cos t i + e sin t j .\]
¡ ¢ ¡ ¢\par
\[\pi\]
Show that the velocity of the particle always makes an angle of with the particle's displace-\par
ment, and that the acceleration of the particle is always perpendicular to its displacement.\par
\[Sketch the path of the particle for 0 6 t 6 \pi.\]
A second particle Q moves on the same path, passing through each point on the path a fixed\par
t\par
time T after P does. Show that the distance between P and Q is proportional to e .\par
\[9465 Jun10 [Turn over\]
\end{flushleft}
\subsection*{Page 5}
\begin{flushleft}
\bigskip\textbf{Question 11.} Two particles of masses m and M , with M > m, lie in a smooth circular groove on a horizontal\par
plane. The coefficient of restitution between the particles is e. The particles are initially\par
projected round the groove with the same speed u but in opposite directions. Find the speeds\par
of the particles after they collide for the first time and show that they will both change\par
\[direction if 2em > M m.\]
-\par
After a further 2n collisions, the speed of the particle of mass m is v and the speed of the\par
particle of mass M is V . Given that at each collision both particles change their directions of\par
motion, explain why\par
\[mv M V = u(M m),\]
- -\par
and find v and V in terms of m, M , e, u and n.\par
\[9465 Jun10\]
\end{flushleft}
\subsection*{Page 6}
\begin{flushleft}
\bigskip\textbf{Section C: Probability and Statistics}\par
\bigskip\textbf{Question 12.} A discrete random variable X takes only positive integer values. Define E(X) for this case,\par
and show that\par
\infty\par
\[E(X) = P (X > n) .\]
\[nX=1\]
I am collecting toy penguins from cereal boxes. Each box contains either one daddy penguin\par
or one mummy penguin. The probability that a given box contains a daddy penguin is p and\par
\[the probability that a given box contains a mummy penguin is q, where p = 0, q = 0 and\]
6 6\par
\[p + q = 1 .\]
Let X be the number of boxes that I need to open to get at least one of each kind of penguin.\par
\[Show that P(X > 4) = p^3 + q^3, and that\]
\[E(X) = 1.\]
pq -\par
\[Hence show that E(X) > 3 .\]
\bigskip\textbf{Question 13.} The number of texts that George receives on his mobile phone can be modelled by a Poisson\par
random variable with mean λ texts per hour. Given that the probability George waits between\par
1 and 2 hours in the morning before he receives his first text is p, show that\par
2λ λ\par
\[pe e + 1 = 0.\]
-\par
Given that 4p < 1, show that there are two positive values of λ that satisfy this equation.\par
The number of texts that Mildred receives on each of her two mobile phones can be modelled\par
by independent Poisson random variables with different means λ 1 and λ 2 texts per hour.\par
Given that, for each phone, the probability that Mildred waits between 1 and 2 hours in the\par
morning before she receives her first text is also p, find an expression for λ 1 + λ 2 in terms of p.\par
Find the probability, in terms of p, that she waits between 1 and 2 hours in the morning to\par
receive her first text.\par
\[9465 Jun10 [Turn over\]
\end{flushleft}
\end{document}
